%!TEX root = main.tex
\setcounter{chapter}{1}
\setcounter{section}{1}
\setcounter{theorem}{0}
\setcounter{equation}{0}


\lectureheader{161}{Calculus I}{Precalculus review}{\textit{Thomas' Calculus} 1.1, 1.2}

\begin{quote}
``To those who do not know mathematics it is difficult to get across a real feeling as to the beauty, the deepest beauty, of nature\dots. 
If you want to learn about nature, to appreciate nature, it is necessary to understand the language that she speaks in."
-- Richard Feynman
\end{quote}

\vfill

\begin{remark}[What is math?]\,
\begin{itemize}
\item Math is \underline{not} just\dots 
\begin{itemize}
\item getting the right answer (though we don't want to be wrong).
\item a magic trick (though some great magic tricks are based on math).
\end{itemize}
\vfill

\item Math is a \underline{language} for building models to describe \underline{precise} relationships.
\begin{itemize}
\item ``Math is the language that with which God has written the universe." -- Galileo
\item Math is the language we use to explain our ideas to other scientists and engineers. 
\begin{itemize}
\item As a undergraduate, Strogatz solved a big open problem concerning DNA by playing with ribbon.
\item His bio-chem professor was excited, but told him he needed a mathematical proof.
\item Published work in \textit{Proceedings of the National Academy of Sciences}.
\end{itemize}
\item Math has its own grammar/syntax.
\item We cannot be precise if we are sloppy.  
\item It does not work to \dots
\begin{itemize}
\item write down most of the correct symbols
	(e.g., ``We calculus.").
\item write down all of the correct symbols but in a strange/nonsensical order
	(e.g., ``to study calculus. going We are").
\end{itemize}
\end{itemize}
\vfill

\item Math is a \underline{science of deduction} that allows us to discover hidden implications of our models.
 \begin{itemize}
 \item Maxwell
 \begin{itemize}
 \item rewrote the experimental results of Faraday and Amp\`ere in terms of calculus.
 \item discovered link between electricity, magnetism, and light.
 \end{itemize}
 \item Gell-man
 \begin{itemize}
 \item modeled known fundamental particles in terms of group theory.
 \item predicted the existence of previously unknown particles.
 \end{itemize}
 \end{itemize}
\vfill
\end{itemize}
\end{remark}

\vfill

\newpage

\begin{remark}[What we need]\,
\begin{itemize}
\item We don't need more humans to do what computers and artificial intelligences can do faster and more reliably.
\item We need more humans that can \underline{understand} the implications of abstract models. 
\item We need more humans who can direct the computers, sift out their garbage responses, and interpret what is meaningful.
\item It's not that you \textit{will} be replaced if you don't adapt -- it's that you \textit{already have} been.
\begin{itemize}
\item We haven't needed human computers since we put the first humans on the moon.
\item We've had software that can perform staggeringly complex calculations for decades now.
\item Now the computers are coming after our ability to ``speak math."
\item Everyday the computers are getting better at stringing together mathematical symbols in syntactically correct combinations -- much more accurately than the average freshman.
\item But the computers still don't know what it \underline{means}.  At least for now, they are merely ``stochastic parrots," and you should want to be more than that.
\end{itemize}

\item And what if the computers do \underline{understand} what they are saying one day?
\begin{itemize}
\item Don't you want to \underline{understand} what the computers figure out about the nature of universe?
\item Should I stop enjoying music just because one AI can identify the song I'm listening to and another can write different song just as pleasing?
\item Does it diminish the enhanced enjoyment that I experience as a result of years of studying music and understanding what I'm listening to on a deeper level?  Surely not.
\item It might just mean that we have a greater supply of teachers and research assistants.
\end{itemize}
\vfill

\end{itemize}
\end{remark}

\newpage

\begin{remark}[What we will do in this class]\,
\begin{itemize}
\item We will learn to write proofs in this class not because I assume that you're all math majors, but because I assume that most of you aren't.
\begin{itemize}
\item STEM professionals use math to \dots
\begin{itemize}
\item write down precise descriptions/models for their objects of interest and the relationships that exist between those objects.
\item deduce hidden consequences of their models or explain experimental consequences already observed.
\item to get  their ideas straight before \underline{translating} them into computer language(s) for large scale calculations.
\end{itemize}
\item In this class, everything is a proof, even a calculation.
\begin{itemize}
\item Write \underline{solutions} as though you are writing a letter to a student in another calculus course at another school who won't believe your result is correct unless you \underline{explain it}.
\item Write/comment your computer code and annotations similarly.
\end{itemize}
\end{itemize}
\vfill

\item Despite the fact that we are not training to be human computers, we will still perform calculations in this class because\dots
\begin{itemize}
\item even though \underline{we don't need} you to be able to do what computers can do faster and more reliably, \underline{you need} to learn how to do those things before you can understand what to ask the computers and how to understand their responses.
\item you need to practice mathematical reasoning in (relatively) concrete settings.
\item you need to develop your ``number sense."
\item you need to develop your ``model sense."
\item I'm deeply suspicious (and you should be too) of claims (even my own) that we no longer need humans who know how to do $X$.
\item new problems are always arising, and old methods have a way of rising from the dead. 
\end{itemize}
\vfill
\end{itemize}
\end{remark}


\newpage

\begin{definition}
A \textbf{(real) function} is a rule $f$ that assigns each member $x$ of set $D$ to a \underline{unique} element $f(x)$ in $\R$.
\begin{itemize}
\item The set of ``input" values $D$ is called the \textbf{domain} of $f$.
\item The set of all ``outputs" $f(D) = \{f(x) : x\in D\}$ is called the \textbf{range} of $f$.
\item We will sometimes write $f:D\to \R$ as a shorthand to say that $f$ is a function with domain $D$.
\item The \textbf{graph} of a function $f:D\to \R$ is the set of ordered pairs
\begin{equation*}
\Big\{\big(x,y\big) : x\in D\text{ and }y=f(x)\Big\}.
\end{equation*}
\end{itemize}
\end{definition}

\begin{example}
Graph the function $y=f(x)$ if
\begin{equation*}
f(x)=\begin{cases}
-x, & x<0\\
x^2,& 0\le x< 1\\
1, & x>1.
\end{cases}
\end{equation*}
\end{example}

\newpage

\begin{theorem}[Vertical line test]
If $f:D\to \R$, then the line $x=a$ intersects the graph of $f$ in exactly one point for each $a$ in $D$.
\end{theorem}

\begin{example}
Sketch the graph of $\big\{(x,y) : x^2+y^2=1\big\}$ and show that it is not a function of $x$.
Finally, cover the graph with two functions of $x$.
\end{example}


\newpage

\begin{definition}
The \textbf{natural domain} of a function $f$ is the largest domain $D$ for which $f(D)\subseteq \R$.
\end{definition}

\begin{example}
Determine the natural domain and associated range for each of the following functions.
\begin{enumerate}
\item $f(x) = x^2$.
\vfill
\item $f(x) = 1/x$.
\vfill
\item $f(x)=\sqrt x$.
\vfill
\item $f(x)=\sqrt{4-x}$.
\vfill
\item $f(x) = \sqrt{1-x^2}$.
\vfill
\end{enumerate}
\end{example}

\newpage

\begin{definition}
Let $f: D\to \R$.
\begin{itemize}
\item We say that $f$ is \textbf{even} if $f(-x) = f(x)$ for all $x\in D$, i.e., the graph of $f$ is symmetric about the $y$-axis.
\item We say that $f$ is \textbf{odd} if $f(-x) = -f(x)$ for all $x\in D$, i.e., the graph of $f$ is symmetric about the origin.
\end{itemize}
\end{definition}


 \begin{example}
Determine (with proof) whether each function is even, odd, or neither.
Hint: You need to compute the domain for each first.
\begin{enumerate}
\item $f(x) = x^2$.
\vfill
\item $f(x)=x^3+x$.
\vfill
\item $f(x) = x+1$.
\vfill
\end{enumerate}
\end{example}

\newpage

\begin{definition}
Let $f:D_f\to \R$ and $g:D_g\to \R$.
\begin{itemize}
\item The rule $(f+g)(x) = f(x) + g(x)$ defines a function with $$\dom(f+g) = D_f\cap D_g.$$
\item The rule $(f-g)(x) = f(x) - g(x)$ defines a function with $$\dom(f-g) = D_f\cap D_g.$$
\item The rule $(fg)(x) = f(x)  g(x)$ defines a function with $$\dom(fg) = D_f\cap D_g.$$
\item The rule $(f/g)(x) = \frac{f(x)}{g(x)}$ defines a function with $$\dom(f/g) = \{x : x\in D_f\cap D_g\text{ and } g(x)\ne 0\}.$$
\end{itemize}
\end{definition}

\begin{example}
Compute $f+g$, $f-g$, $fg$ and $f/g$ and their domains if $f(x)=\sqrt x$ and $g(x)=\sqrt{1-x}$.
\end{example}

\newpage

\begin{definition}
Let $f:D_f\to \R$ and $g:D_g\to \R$.
The rule $(f\circ g)(x) = f(g(x))$ defines a function with
$$\dom(f\circ g) = \{x: x\in D_g\text{ and } g(x)\in D_f\}.$$
\end{definition}

\begin{example}
Compute $f\circ g$ together with its domain for the following.
\begin{enumerate}
\item $f(x) =x^2+1$ and $g(x)=\DS\frac{x+1}{x-1}$.
\vfill
\item $f(x)=x^2$ and $g(x)=\sqrt x$.
\vfill
\item $f(x)=\sqrt x$ and $g(x)=x^2$.
\vfill
\end{enumerate}
\end{example}

\newpage

\begin{example}
Given the function $f(x)=x^4-4x^3+10$, find formulas to \dots
\begin{enumerate}
\item compress the graph horizontally by a factor of 2 followed by a reflection across the $y$-axis and then a shift 3 units to the right.
\vfill
\item compress the graph vertically by a factor of 2 followed by a reflection across the $x$-axis and then a shift 1 unit up.
\vfill
\end{enumerate}
\end{example}

